%! TEX program = lualatex
%! TEX root = thesis.tex

% DocumentMetadata is required for an accessible PDF
\DocumentMetadata{
    lang=en,
    pdfstandard=ua-2,
    tagging=on,
    tagging-setup={
        math/setup={mathml-AF,mathml-SE},
        extra-modules={verbatim-mo},
    }
}
\documentclass{uofithesis}
% \documentclass[thesis]{uofithesis} % for a thesis

\usepackage{fontsetup}

% If using math unicode-math is required for math accessibility.
% If not using math, remove this package.
\usepackage{unicode-math}

% Remove any graphics packages you are not using.
% Graphics packages can be used as long as alt-text is provided.
% For tikz alt text can be added with:
%   \begin{tikzpicture}[alt={Alt text...}]
% For images:
%  \includegraphics[alt={Alt text...}]{image.png}
\usepackage{graphicx}
\usepackage{tikz}
% uofithesis sets accessible colors for use in charts, but you can also define 
% your own if needed
\usepackage{pgfplots}

\usepackage{csquotes}

% Bibliography for references
% biber must be set as the backend, style can be changed as needed
\usepackage[hidelinks,pdfusetitle]{hyperref}
\usepackage[backend=biber,style=ieee]{biblatex}
\usepackage[bottom]{footmisc}

% At time of writing, chemistry packages are not compatible with accessibility 
% tagging, but there are no alternatives.
\usepackage[version=4]{mhchem}
\usepackage{chemfig}
\pgfplotsset{compat=1.18}

\addbibresource{./references.bib}

% Raggedright is recommended for accessibility
\raggedright
\setlength{\parindent}{1.5em}

\title{Coffee Consumption Among Graduate Students Writing Dissertations}
\author{Alice Ascher}
\department{Food Science and Human Nutrition}

% You can set a specific earned degree instead of Master/Doctor of Science/Philosophy
%\earneddegree{Doctor of Education}

% Add a concentration if applicable
%\concentration{Coffee Studies}
% Add a minor if applicable
%\minor{Global Studies}

\degreeyear{2026}

% You can set a custom text for the committee section, otherwise it will default 
% to "Doctoral Committee:" or "Master's Committee:" based on the degree level
%\committeetitle{Advisers:}

\committee{
\item Associate Professor Bob Barnard, Chair
\item Assistant Professor Charlie Clarke
\item Professor Dan Drower
\item Professor Eve Earlsfield
}

\begin{document}

% Uncomment the copyrightpage macro to include an automatic copyright page
% before the title
%\copyrightpage
\maketitle

\abstract{
    This is the abstract of the thesis. It should be a concise summary of the 
    research, including the main findings and conclusions. The abstract should 
    be no more than 350 words in length.
}

\acknowledgements{
    This page is optional. The text should be in either 10- or 12-point font 
    and the line spacing should be either 1.5 or double-spaced.

    If you do not want to include an acknowledgment, delete the text of the 
    acknowledgment and the \texttt{\textbackslash acknowledgements} command 
    from your \texttt{.tex} file.
}

\tableofcontents


% Sections make up the Main Text. Always use sectioning commands to create sections, 
% subsections, etc. Do not use formatting commands to make regular text look like 
% a section heading, as this will not be tagged properly for accessibility.
\section{Coffee Consumption Mathematics}

Coffee extraction occurs when hot water is poured over coffee grounds, causing 
desirable compounds such as caffeine, carbohydrates, lipids, melanoidins and 
acids to be extracted from the grounds.

Extraction yield refers to the solubles dissolved during brewing. This is often 
expressed as a percentage of the coffee's mass. It is also known as solubles 
yield or simply extraction. The extraction yield percentage describes the mass 
transferred from coffee grounds to water, expressed as a percentage of the 
initial mass of the grounds. It is given by (\ref{eq:extraction_yield}).

\tagstructbegin{tag=Part}
\begin{equation}\label{eq:extraction_yield}
    Y = \frac{M_2 \cdot t}{M_1}
\end{equation}
\tagstructend

where $Y$ is the extraction yield expressed as a percentage, $t$ is the total 
dissolved solids expressed as a percentage of the final beverage, $M_1$ is the 
mass of the grounds in grams, and $M_2$ is the final beverage's mass in grams. 
This means that an extraction yield of 20\% can be obtained by brewing 18 grams 
of coffee, resulting a 36-gram final beverage with a $t$ of 10\%. Yield can also 
be expressed as total dissolved solids, or parts-per-million (ppm).

\section{Coffee Berries}

A coffee bean is a seed from the Coffea plant and the source for coffee. This 
fruit is often referred to as a coffee cherry, but unlike the cherry, which 
usually contains a single pit, it is a berry most commonly found having two 
seeds with their flat sides together.

\tagstructbegin{tag=Part}
\begin{figure}[H]
    \centering
    \includegraphics[width=0.6\textwidth,alt={Multiple red and green coffee 
    berries clustered together on coffee tree branches}]{
        Clumps_of_coffee_berries_growing_on_a_tree_by_Brian_Smith.jpg
        }
    \caption{Clumps of coffee berries growing on a tree (Photo by Brian Smith)}
    \label{fig:coffee_berries}
\end{figure}
\tagstructend

\subsection{Characteristics of Coffee Berries}

The two most economically important varieties of coffee plants are the arabica 
and the robusta; approximately 60\% of the coffee produced worldwide is arabica 
and some 40\% is robusta.

\tagstructbegin{tag=Part}
\begin{table}[H]
    \centering
    \caption{Common Coffee Species Comparison}
    \medskip
    \label{tab:coffee_species}
    \renewcommand{\arraystretch}{1.5}
    \tagpdfsetup{table/header-columns={1},table/header-rows={1}}
    \begin{tabular}{|>{\bfseries}l|l|l|}
        \hline
        Species & \textbf{Flavor Profile} & \textbf{Global Production} \\
        \hline
        Arabica & Smooth, complex, with fruit notes & 60\% \\
        \hline
        Robusta & Strong, bitter, earthy & 40\% \\
        \hline
    \end{tabular}
\end{table}
\tagstructend
\pagebreak

The following table summarizes key information about coffee berries:

\tagstructbegin{tag=Part}
\begin{table}[H]
    \centering
    \caption{Key Characteristics of Coffee Berries}
    \medskip
    \label{tab:coffee_characteristics}
    \renewcommand{\arraystretch}{1.5}
    \tagpdfsetup{table/header-columns={1}}
    \begin{tabular}{|>{\bfseries}p{2in}|p{4in}|}
        \hline
        Color at Ripeness & Bright red, yellow or pink depending on variety \\
        \hline
        Size & Approximately 15 mm in diameter \\
        \hline
        Weight & 300--330 mg dry \\
        \hline
        Seeds per Berry & Typically 2 seeds; 10--15\% are peaberries (single seed) \\
        \hline
        Time to Maturity & 6--11 months depending on variety and growing conditions \\
        \hline
        Caffeine Content & 1.0--2.5\% by dry weight \\
        \hline
        Growing Conditions & Altitudes of 600--2,200 meters; temperatures of 15--24°C \\
        \hline
        Shelf Life (Fruit) & 2--3 weeks if left on tree after ripening \\
        \hline
    \end{tabular}
\end{table}
\tagstructend

\section{Coffee Chemistry}

Caffeine \ce{C8H10N4O2} is a central nervous system (CNS) stimulant of the 
methylxanthine class and is the most commonly consumed psychoactive substance 
globally.

\tagstructbegin{tag=Part}
% molecule from https://tug.ctan.org/macros/generic/chemfig/chemfig-en.pdf
\begin{figure}[H]
    \centering
    \chemfig{*6((=O)-N(-CH_3)-*5(-N=-N(-CH_3)-=)--(=O)-N(-H_3C)-)}
    \caption{Molecular structure of caffeine (\ce{C8H10N4O2})}
    \label{fig:caffeine_molecule}
\end{figure}
\tagstructend

Some people
\footnote{Over 90\% of American adults consume caffeine daily
\cite{SleepFoundation1}.}
drink beverages containing caffeine to relieve or prevent drowsiness and to 
improve cognitive performance. To make these drinks, caffeine is extracted by 
steeping the plant product in water, a process called infusion.

Overall population mean coffee consumption was stable from 2003 to 2012
\cite{LOFTFIELD20161762}.

\section{Coffee Charts}

Data visualization is an important aspect of understanding coffee consumption 
patterns and production statistics. The following chart displays typical daily 
coffee consumption at different times.

\tagstructbegin{tag=Part}
\begin{figure}[H]
    \begin{tikzpicture}[alt={Bar chart showing a marked increase in coffee 
        consumption during deposit week with a steep decline after deposit}]
        \begin{axis}[
                width=0.9\textwidth,
                height=3in,
                ybar,
                cycle list name=StrictAA,
                every axis plot/.append style={fill, draw=black},
                nodes near coords style={text=black, font=\small},
                enlargelimits=0.15,
                ylabel={Avg. Cups per Day},
                symbolic x coords={Weekday, Weekend, Deposit Week, Post-Deposit},
                xtick=data,
                nodes near coords,
                nodes near coords align={vertical},
                bar width=20pt
            ]
            \addplot+ coordinates {(Weekday,2.1) (Weekend,1.4) (Deposit Week,4.6) (Post-Deposit,0.3)};
        \end{axis}
    \end{tikzpicture}
    \caption{Average Daily Coffee Consumption By Period}
    \medskip
    \label{fig:coffee_consumption_chart}
\end{figure}
\tagstructend
\bigskip

The following line chart illustrates the trend of coffee consumption throughout 
a week. There is a noticeable increase in consumption during the middle of the 
week, peaking on Thursday and Friday, followed by a decline over the weekend.

\tagstructbegin{tag=Part}
\begin{figure}[H]
    \centering
    \begin{tikzpicture}[alt={Line chart showing the trend of coffee consumption 
        throughout the week, with a peak on Thursday and Friday}]
        \begin{axis}[
                width=0.9\textwidth,
                height=5cm,
                xlabel={Day of Week},
                ylabel={Cups Consumed},
                legend style={at={(0.5,-0.15)},
                anchor=north,legend columns=-1},
                grid=major,
                xtick={1,2,3,4,5,6,7},
                xticklabels={Mon, Tue, Wed, Thu, Fri, Sat, Sun},
            ]
            \addplot[mark=*] coordinates {(1,2.4) (2,2.6) (3,2.8) (4,3.2) (5,3.5) (6,1.8) (7,1.2)};
            \legend{Daily Intake}
        \end{axis}
    \end{tikzpicture}
    \caption{Coffee Consumption Trend Throughout the Week}
    \medskip
    \label{fig:coffee_trend_chart}
\end{figure}
\tagstructend

\printbibliography[heading=bibintoc,title={References}]

% For an appendix without the chapter numbering, use \section* and add it to 
% the table of contents manually with \addcontentsline
\section*{Appendix A: Coffee Survey}
\addcontentsline{toc}{section}{Appendix A: Coffee Survey}

\begin{enumerate}
    \item How many cups of coffee do you consume on an average weekday?
    \item How many cups of coffee do you consume on an average weekend day?
    \item How many cups of coffee do you consume during deposit week?
    \item How many cups of coffee do you consume after deposit week?
\end{enumerate}

\end{document}